\documentclass{article}

%%%%%%%%%%%%%%%%%%%%%%%%%%%%%%%%%%%%%%%%%%%%%%%%%%%%%%%%%%%%%%%%%%%%%%%%%%%%%%%%%%%%%%%%
% RPGCOLOR PACKAGE 
%%%%%%%%%%%%%%%%%%%%%%%%%%%%%%%%%%%%%%%%%%%%%%%%%%%%%%%%%%%%%%%%%%%%%%%%%%%%%%%%%%%%%%%%

% Load the package. Without an option, the package selects the default theme.
\usepackage{rpgcolor}

% You can also choose a built-in theme as a package option, for example:
% \usepackage[theme=PHB2014]{rpgcolor}
% \usepackage[PHB2014]{rpgcolor}

% Create a document-local theme with the same semantic color slots used by the
% built-in themes.

\definecolor{MyCoolColor}{HTML}{DADB0D} % RPGColor already loads xcolor
\RPGColorThemeCreate{custom}
{%
    base      = MyCoolColor, % You can use custom colors
    comment   = RPGColorLavender,
    sidebar   = RPGColorCeladon,
    tablerow  = RPGColorOlive,
    narration = RPGColorCoral,
    title     = RPGColorAubergine,
    hrule     = RPGColorDeepTeal,
    statblock = RPGColorFrost,
    ribbon    = RPGColorAmber,
    contour   = RPGColorAsh,
    footer    = RPGColorOxblood,
}

%%%%%%%%%%%%%%%%%%%%%%%%%%%%%%%%%%%%%%%%%%%%%%%%%%%%%%%%%%%%%%%%%%%%%%%%%%%%%%%%%%%%%%%%
%   DEMO HELPERS
%%%%%%%%%%%%%%%%%%%%%%%%%%%%%%%%%%%%%%%%%%%%%%%%%%%%%%%%%%%%%%%%%%%%%%%%%%%%%%%%%%%%%%%%

% Take a generic xcolor and make a row with | \ColorName | <Color Box> |
\newcommand{\ColorRow}[1]
{%
    \rule[-0.35em]{0pt}{1.5em}\texttt{\char`\\#1}
    & \colorbox{#1}{\rule[-0.35em]{0pt}{1.5em}\hspace{18em}}\\
}

% Use \ColorRow to demo a whole RPGColor theme in a table
\newcommand{\ThemeTable}[1]
{%
    \RPGColorThemeApply{#1}%
    \subsection*{\texttt{#1}}

    \begin{tabular}{|l|l|}
        \hline
        Color control sequence & Color \\
        \hline
        \ColorRow{RPGColorBase}
        \ColorRow{RPGColorComment}
        \ColorRow{RPGColorSidebar}
        \ColorRow{RPGColorTableRow}
        \ColorRow{RPGColorNarration}
        \ColorRow{RPGColorTitle}
        \ColorRow{RPGColorHRule}
        \ColorRow{RPGColorStatBlock}
        \ColorRow{RPGColorRibbon}
        \ColorRow{RPGColorContour}
        \ColorRow{RPGColorFooter}
        \hline
    \end{tabular}
}

%%%%%%%%%%%%%%%%%%%%%%%%%%%%%%%%%%%%%%%%%%%%%%%%%%%%%%%%%%%%%%%%%%%%%%%%%%%%%%%%%%%%%%%%

\begin{document}

\section*{RPG Color Themes}

Each table switches theme with the \texttt{\char`\\RPGColorThemeApply} command and then displays the semantic color control sequences defined by that theme. See the README for a full list of themes.

\ThemeTable{default}
\ThemeTable{PHB2014}
\ThemeTable{DMG2014}
\ThemeTable{arcane}
\ThemeTable{feywild}
\ThemeTable{infernal}
\ThemeTable{maritime}
\ThemeTable{custom}

\subsection*{Local Overrides}

Use \texttt{\char`\\RPGColorSet} with \texttt{key=value,...} when you want to change one or more live semantic colors without changing the saved theme definition.

\RPGColorThemeApply{default}%
\RPGColorSet{title=RPGColorForest,footer=RPGColorInk,ribbon=MyCoolColor}%

This table starts from \texttt{default}, then changes only \texttt{title}, \texttt{footer}, and \texttt{ribbon} with \texttt{\char`\\RPGColorSet}.
The saved theme is unchanged; applying it again restores the original values.

\begin{tabular}{|l|l|}
    \hline
    Color control sequence & Color \\
    \hline
    \ColorRow{RPGColorBase}
    \ColorRow{RPGColorComment}
    \ColorRow{RPGColorSidebar}
    \ColorRow{RPGColorTableRow}
    \ColorRow{RPGColorNarration}
    \ColorRow{RPGColorTitle}
    \ColorRow{RPGColorHRule}
    \ColorRow{RPGColorStatBlock}
    \ColorRow{RPGColorRibbon}
    \ColorRow{RPGColorContour}
    \ColorRow{RPGColorFooter}
    \hline
\end{tabular}

\end{document}
